\scanpage{3}

Une façon de le voir est de dire que
la donnée de $\uline{s}$ définit un scindage
\inserted{$\varphi_{\uline{s}}$} de
l'extension $(*)$ ci-dessus, et que \unsure{vu pour}
ss-groupe de $\operatorname{SRT}(X)$,
$S_{\unsure{1|*}}T(X,\uline{s})$ s'identifie
à l'image inverse de
\[
\varphi_{\uline{s}}(T(X))\subset\operatorname{SUT}(X)
\]
par l'hom. can
\[
\operatorname{SUR}(X)\longrightarrow\operatorname{SUT}(X)
\]
de passage au quotient. Le chgt de
$\uline{s}$ en $\uline{s}'=u.\uline{s}$
transforme $\varphi_{\uline{s}}$ en
\[
\varphi_{\uline{s}'}=\operatorname{int}(u)\circ\varphi_{\uline{s}}
\]
donc si $\widetilde u\in\mathbb R^I$ relève $u$, on trouve
\[
S_{\unsure{1|*}}T(X,\uline{s}')
=\operatorname{int}(\widetilde u)S_{\unsure{1|*}}T(X,\uline{s}).
\]

En termes de $\vec X$, la définition de
$S_{\unsure{1|*}}T(X,\uline{s})$
dépend d'un choix $\uline{s}$ \uline{non-anodin},
par contre celle de $S_1RT(X)$ ne dépend que d'
un choix anodin, celui de la famille des
structures toriques $(t_i)_{i\in I}$.
De plus $\prod T_i$ comme
torseur sous $\unsure{\mathbb U^I}$ se réalise
\inserted{(pas fidèlement d'ailleurs)}
comme une classe de scindages de l'ext $(*)$
(via autom. int. par $\mathbb U^I$)
$\bigl[$\unsure{déduits} de $\operatorname{SRT}(X)$
avec \unsure{ses} structures suppl.$\bigr\}$
savoir les inclusions

\diagram{
\begin{tikzcd}[row sep=small]
\mathbb R^I\arrow[dr] &\\
&\operatorname{SRT}(X)\\
S_{\unsure{1|*}}T^!(X)\arrow[ur]&
\end{tikzcd}
}

\unsure{coïncid.} sur $\mathbb Z^I$.

\margin{2}
